\chapter{Conclusion and Future Scope}\label{conclusion}
The proposed adaptive insider threat detection framework combines causal analysis, MAD-based feature refinement, and LNNs to a large extent to solve the five major issues raised by the insider threat detection problem, namely class imbalance, temporal variability, interpretability, scalability, and generalization to the new users problem. The experiments performed at the granularity of session, day, and behavior show consistent and balanced improvements of precision, recall, and F1-score with a low FPR. In addition, the profiling for runtime and memory indicates that the proposed model is efficient in terms of the required computational resources and the speed of convergence as compared to the traditional RNN and LSTM models. 

By combining causal reasoning with MAD-based refinement, the system becomes more interpretable as it focuses on the most statistically influential behavioral attributes, while the transition to an adaptive LNN backbone enables scalable and temporally adaptable behavior. The components, taken together, empower the framework to effectively detect subtle insider behaviors across diverse operational scenarios. 
While these successes are notable, the authors still acknowledge the existence of several limitations. The evaluation is heavily dependent on the CMU-CERT synthetic datasets, which, even though they are the standard benchmark in this research area, still cannot adequately simulate the complexity and heterogeneity of the real-world enterprise environments. Furthermore, the interpretability achieved through causal-MAD reasoning is currently only statistical and is yet to be developed into semantically rich, analyst-friendly explanations. Lastly, the framework’s ability to adapt to the changing behavioral patterns and the evolving tactics of insiders in the long run is still a challenge that has not been addressed. 

The work ahead will concentrate on two major aspects: Firstly, the methodology proposed will be tested on real-world, multimodal enterprise datasets that combine HR, access control, and network traffic logs to guarantee higher operational robustness. Secondly, causal-semantic explanation components will be implemented to link statistical insights with domain knowledge, thus facilitating human–AI collaborative decision-making. The aggregate effect of these changes is a framework that is closer to being a viable, scalable, and explainable solution for the enterprise-level detection of insider threats which ensures security monitoring that is proactive, transparent, and trustworthy.
